\documentclass[aspectratio=169, 10pt]{beamer}

% --- Paquetes Fundamentales ---
\usepackage[utf8]{inputenc}
\usepackage[T1]{fontenc}
\usepackage[spanish,es-tabla]{babel}
\usepackage{amsmath, amssymb}
\usepackage{microtype}
\usepackage{booktabs}

% --- Matemáticas y Electrónica ---
\usepackage{siunitx}
\usepackage[american, RPvoltages]{circuitikz}

% --- Código Fuente ---
\usepackage{listings}
\usepackage{xcolor}
\usepackage{colortbl}

% --- Configuración de siunitx ---
\sisetup{output-decimal-marker = {,}, per-mode = symbol}

% --- Configuración de Colores e Identidad ---
\definecolor{ucbblue}{RGB}{0, 51, 102}
\definecolor{codegreen}{rgb}{0,0.6,0}
\definecolor{codegray}{rgb}{0.5,0.5,0.5}
\definecolor{codepurple}{rgb}{0.58,0,0.82}

\usetheme{Boadilla}
\usecolortheme{seahorse}
\setbeamertemplate{navigation symbols}{}
\setbeamertemplate{itemize items}[circle]
\setbeamercolor{title}{bg=ucbblue, fg=white}
\setbeamercolor{frametitle}{bg=ucbblue, fg=white}

% --- Macros UI ---
\newcommand{\tecla}[1]{\colorbox{gray!20}{\texttt{\footnotesize\textbf{#1}}}}

% --- Configuración de Listings (Python) ---
\lstset{
    language=Python,
    basicstyle=\ttfamily\scriptsize,
    keywordstyle=\color{blue}\bfseries,
    commentstyle=\color{codegreen}\itshape,
    stringstyle=\color{codepurple},
    numberstyle=\tiny\color{codegray},
    numbers=left,
    stepnumber=1,
    showstringspaces=false,
    frame=single,
    backgroundcolor=\color{gray!5},
    rulecolor=\color{gray!40},
    breaklines=true,
    aboveskip=2pt,
    belowskip=2pt
}

% --- Metadatos ---
\title[Análisis Nodal Modificado (MNA)]{Análisis Nodal Modificado (MNA) y Redes Activas Lineales}
\subtitle{Supernodos, Fuentes Dependientes y Caracterización de Transferencia}
\author[M.Sc. Ing. B. Quiroga]{M.Sc. Ing. Bernardo Quiroga Turdera}
\institute[UCB Tarija]{IMT-121 Circuitos Electrónicos I | UCB Sede Tarija}
\date{Sesión 07}

\begin{document}

% ==========================================
% SLIDE 1: Portada
% ==========================================
\begin{frame}
    \titlepage
\end{frame}

% ==========================================
% SLIDE 1b: Resultados de Aprendizaje y Ruta
% ==========================================
\begin{frame}{Resultados de Aprendizaje y Alineamiento}
    \textbf{Alineamiento Curricular:} Criterio ABET 1 (Modelado Matemático Complejo) y Marco CDIO (Concebir–Diseñar).
    
    \begin{block}{Resultado de Aprendizaje (RA)}
        Formular sistemas de ecuaciones nodales modificadas para redes con fuentes dependientes y supernodos, implementando el solver computacional en Python y la directiva \texttt{.tf} en LTspice para extraer funciones de transferencia e impedancias.
    \end{block}

    \vspace{0.2cm}
    \begin{columns}
        \begin{column}{0.9\textwidth}
            \textbf{Hoja de Ruta de la Sesión:}
            \begin{enumerate}
                \item ¿Por qué colapsa la Inspección Directa clásica?
                \item Tratamiento riguroso de Supernodos (Conservación + Restricción).
                \item Modelado de Fuentes Dependientes (Preamplificadores y Sensores).
                \item Resolución analítica en pizarra ($2\times 2$ aumentado).
                \item Co-Programación en Python (\texttt{numpy.linalg.solve}).
                \item Taller de LTspice: Extracción de $A_v, R_{\text{in}}, R_{\text{out}}$ (\texttt{.tf}).
            \end{enumerate}
        \end{column}
    \end{columns}
\end{frame}

% ==========================================
% SLIDE 2: Limitaciones de la Inspección Directa
% ==========================================
\begin{frame}{De la Inspección Directa al MNA}
    \textbf{El Límite de la Inspección Directa ($\mathbf{G}\mathbf{V} = \mathbf{I}$):}
    \begin{itemize}
        \item La formulación asume que todas las ramas siguen la Ley de Ohm ($I = G \cdot \Delta V$).
        \item Si una fuente de tensión (independiente o controlada) no tiene resistencia en serie, su impedancia es $R \to 0$ ($G \to \infty$), provocando una \textbf{indeterminación matemática}.
    \end{itemize}

    \vspace{0.3cm}
    \begin{table}[]
        \centering
        \footnotesize
        \renewcommand{\arraystretch}{1.4}
        \begin{tabular}{@{}p{3.8cm}p{4cm}p{5.2cm}@{}}
            \toprule
            \rowcolor{ucbblue!10}
            \textbf{Escenario Topológico} & \textbf{Comportamiento en Matriz G} & \textbf{Solución Metodológica (MNA)} \\
            \midrule
            Fuente a Tierra (GND) & Fija directamente el voltaje nodal. & Reduce la dimensión ($V_{\text{nodo}} = V_s$). \\
            Fuente entre 2 nodos flotantes & Corriente indefinida por Ohm. & Formación de \textbf{Supernodo} + Ec. de Enlace. \\
            Fuentes dependientes & Rompe la simetría ($G_{ij} \neq G_{ji}$). & Inclusión de \textbf{variables de control} en matriz $\mathbf{A}$. \\
            \bottomrule
        \end{tabular}
    \end{table}
\end{frame}

% ==========================================
% SLIDE 3: Supernodos
% ==========================================
\begin{frame}{Tratamiento Físico-Matemático de Supernodos}
    \textbf{Definición de Supernodo:}
    Superficie cerrada que engloba a una fuente de voltaje conectada entre dos nodos no referenciados, junto con cualquier elemento en paralelo.

    \vspace{0.2cm}
    \begin{center}
        \begin{circuitikz}[scale=0.8, transform shape]
            \draw (0,0) node[left]{Nodo 2} to[short, o-] (1.5,0)
                  to[V, v=$V_{\text{fuente}}$] (4.5,0)
                  to[short, -o] (6,0) node[right]{Nodo 3};
            % Caja del supernodo
            \draw[dashed, red, thick, rounded corners] (0.5, -1) rectangle (5.5, 1);
            \node[red, above] at (3, 1) {\textbf{SUPERNODO (2-3)}};
        \end{circuitikz}
    \end{center}

    \vspace{0.1cm}
    \begin{block}{Las Dos Ecuaciones Obligatorias del Supernodo}
        \small
        \begin{enumerate}
            \item \textbf{Ecuación de Conservación (LCK):} Suma de corrientes que \textit{escapan} del supernodo: 
            $$ \sum I_{\text{salientes (Nodo 2)}} + \sum I_{\text{salientes (Nodo 3)}} = 0 $$
            \item \textbf{Ecuación de Restricción / Enlace:} Relación impuesta por la fuente:
            $$ V_{\text{positivo}} - V_{\text{negativo}} = V_{\text{fuente}} $$
        \end{enumerate}
    \end{block}
\end{frame}

% ==========================================
% SLIDE 4: Fuentes Dependientes (Refactorizado)
% ==========================================
\begin{frame}{Fuentes Dependientes como Modelos de Transductores}
    En mecatrónica y biomédica, los transductores activos y etapas de amplificación se modelan mediante fuentes controladas lineales:

    \vspace{0.4cm}
    \begin{columns}
        % Columna Izquierda (Fuentes de Voltaje)
        \begin{column}{0.5\textwidth}
            \centering
            \textbf{VCVS} (Tensión / Tensión)\\
            \begin{circuitikz}[scale=0.7, transform shape]
                \draw (0,0) to[cV, v=$A_v V_x$] (3,0); 
            \end{circuitikz} \\
            \textcolor{blue}{\footnotesize (Control: $V_x$)}
            
            \vspace{0.6cm}
            \textbf{CCVS} (Tensión / Corriente)\\
            \begin{circuitikz}[scale=0.7, transform shape]
                \draw (0,0) to[cV, v=$R_m I_x$] (3,0); 
            \end{circuitikz} \\
            \textcolor{blue}{\footnotesize (Control: $I_x$)}
        \end{column}
        
        % Columna Derecha (Fuentes de Corriente)
        \begin{column}{0.5\textwidth}
            \centering
            \textbf{VCCS} (Corriente / Tensión)\\
            \begin{circuitikz}[scale=0.7, transform shape]
                \draw (0,0) to[cI, i=$g_m V_x$] (3,0); 
            \end{circuitikz} \\
            \textcolor{blue}{\footnotesize (Control: $V_x$)}
            
            \vspace{0.6cm}
            \textbf{CCCS} (Corriente / Corriente)\\
            \begin{circuitikz}[scale=0.7, transform shape]
                \draw (0,0) to[cI, i=$\beta I_x$] (3,0); 
            \end{circuitikz} \\
            \textcolor{blue}{\footnotesize (Control: $I_x$)}
        \end{column}
    \end{columns}

    \vspace{0.5cm}
    \small \textbf{Regla Matricial:} Identificar la variable de control, expresarla en función de los potenciales nodales, y reagrupar términos a la izquierda para mantener $\mathbf{A}\mathbf{x} = \mathbf{b}$.
\end{frame}

% ==========================================
% SLIDE 5a: Ejemplo Integrador (Planteamiento)
% ==========================================
\begin{frame}{Ejercicio Integrador: Planteamiento y Restricción}
    \textbf{Circuito:} Preamplificador con VCVS ($E_1 = 3V_x$). Señal de control $V_x = V_2$.

    \vspace{0.2cm}
    \begin{columns}
        \begin{column}{0.45\textwidth}
            \begin{center}
                \begin{circuitikz}[scale=0.7, transform shape]
                    \draw (0,0) to[V, v=$V_s {=} \qty{10}{\volt}$] (0,3) node[above]{Nodo 1}
                          to[R, l=$R_1 {=} \qty{50}{\ohm}$] (3,3) node[above]{Nodo 2}
                          to[cV, v=$3V_x$] (6,3) node[above]{Nodo 3}
                          to[R, l=$R_3 {=} \qty{200}{\ohm}$] (6,0) -- (0,0);
                    \draw (3,3) to[R, l=$R_2 {=} \qty{100}{\ohm}$] (3,0);
                    \draw (0,0) node[ground]{};
                    \draw (3,0) node[ground]{};
                    \draw (6,0) node[ground]{};
                    \draw[<->, dashed, blue] (3.4, 2.8) -- (3.4, 0.2) node[midway, right]{$V_x$};
                \end{circuitikz}
            \end{center}
        \end{column}
        \begin{column}{0.55\textwidth}
            \small
            \textbf{Deducción Analítica (Parte 1):}
            \begin{itemize}
                \item \textbf{Nodo 1:} Fijo por fuente $\implies V_1 = \qty{10.0}{\volt}$.
                \item \textbf{Variable de Control:} $V_x = V_2$.
            \end{itemize}
            
            \vspace{0.2cm}
            \textbf{Ecuación de Enlace del Supernodo (2-3):}
            \begin{align*}
                V_2 - V_3 &= 3 V_x \\
                V_2 - V_3 &= 3 V_2 \\
                -2 V_2 - V_3 &= 0 \quad \text{\textbf{--- [Ecuación 1]}}
            \end{align*}
        \end{column}
    \end{columns}
\end{frame}

% ==========================================
% SLIDE 5b: Ejemplo Integrador (Resolución)
% ==========================================
\begin{frame}{Ejercicio Integrador: Conservación LCK y Solución}
    \textbf{Deducción Analítica (Parte 2):}
    
    LCK en la frontera del Supernodo (2-3). Sumamos las corrientes que salen del Nodo 2 y del Nodo 3:
    \begin{equation*}
        \frac{V_2 - 10}{50} + \frac{V_2}{100} + \frac{V_3}{200} = 0
    \end{equation*}
    
    Multiplicando todo por $200$ para eliminar denominadores algebraicos:
    \begin{align*}
        4(V_2 - 10) + 2 V_2 + V_3 &= 0 \\
        6 V_2 + V_3 &= 40 \quad \text{\textbf{--- [Ecuación 2]}}
    \end{align*}
    
    \vspace{0.1cm}
    \begin{alertblock}{Sistema Matricial Resultante}
        $$
        \begin{bmatrix} 6.0 & 1.0 \\ -2.0 & -1.0 \end{bmatrix} \begin{bmatrix} V_2 \\ V_3 \end{bmatrix} = \begin{bmatrix} 40.0 \\ 0.0 \end{bmatrix} \implies \begin{cases} \mathbf{V_2 = +10.00\text{ V}} \\ \mathbf{V_3 = -20.00\text{ V}} \end{cases}
        $$
    \end{alertblock}
\end{frame}

% ==========================================
% SLIDE 6: Python (MNA)
% ==========================================
\begin{frame}[fragile]{Implementación Computacional en Python (NumPy)}
\begin{lstlisting}
import numpy as np

# 1. Definicion del Sistema A * V = b (Variables: V2, V3)
# Fila 1 (LCK Supernodo):     6*V2 +  1*V3 = 40
# Fila 2 (Restriccion VCVS): -2*V2 -  1*V3 = 0
A = np.array([[ 6.0,  1.0], [-2.0, -1.0]], dtype=float)
b = np.array([40.0, 0.0], dtype=float)

# 2. Verificacion de singularidad y solucion
det_A = np.linalg.det(A)
assert np.abs(det_A) > 1e-6, "Error: Sistema singular (det = 0)."
V = np.linalg.solve(A, b)

# 3. Metricas de ingenieria
Av = V[1] / 10.0  # Ganancia de voltaje respecto a Vs (10V)

print("--- RESULTADOS COMPUTACIONALES (MNA) ---")
print(f"Determinante: {det_A:.2f}")
print(f"Voltaje V2 (Control Vx): {V[0]:.3f} V")
print(f"Voltaje V3 (Salida Vo) : {V[1]:.3f} V")
print(f"Ganancia del Sistema Av: {Av:.3f} V/V")
\end{lstlisting}
\end{frame}

% ==========================================
% SLIDE 7: LTspice .tf
% ==========================================
\begin{frame}{Taller de LTspice — Directiva \texttt{.tf} (Transfer Function)}
    La directiva \texttt{.tf} evalúa la ganancia en pequeña señal, impedancia de entrada y resistencia de salida de forma completamente automatizada.

    \vspace{0.2cm}
    \begin{center}
        \begin{circuitikz}[scale=0.75, transform shape]
            % Salida del VCVS
            \draw (2,2) node[left]{(Pin +)} to[short, o-] (3,2) 
                  to[cV, v=e (\texttt{VCVS})] (6,2) 
                  to[short, -o] (7,2) node[right]{(Pin -)};
            \node[above] at (2,2.2) {Nodo 2}; \node[above] at (7,2.2) {Nodo 3};
            
            % Control del VCVS
            \draw (2,0) node[left]{(Ctrl +)} to[short, o-o] (7,0) node[right]{(Ctrl -)};
            \node[below] at (2,-0.2) {Nodo 2}; \node[below] at (7,-0.2) {GND};
            
            \draw[dashed, gray] (4.5, 0.2) -- (4.5, 1.2); % Indicador de control
        \end{circuitikz}
    \end{center}

    \vspace{0.1cm}
    \begin{columns}
        \begin{column}{0.4\textwidth}
            \small
            \textbf{Protocolo de Simulación:}
            \begin{enumerate}
                \item Insertar el componente \texttt{e} (\tecla{F2}).
                \item Asignar \texttt{Value = 3}.
                \item Directiva: \texttt{.tf V(N003) V1}
            \end{enumerate}
        \end{column}
        \begin{column}{0.6\textwidth}
            \small
            \textbf{Interpretación del Reporte SPICE:}
            \begin{itemize}
                \item \texttt{Transfer\_function = -2.000} ($V_3/V_1 = -2$).
                \item \texttt{v1\#Input\_impedance = 83.333 ohms}.
                \item \texttt{output\_impedance\_at\_v(n003) = 0 ohms} (fuente ideal).
            \end{itemize}
        \end{column}
    \end{columns}
\end{frame}

% ==========================================
% SLIDE 8: Dinámica en Aula
% ==========================================
\begin{frame}{Dinámica Activa en Aula (Tríadas PFI)}
    \small
    \textbf{Desafío Paramétrico (15 minutos):}
    Un transductor sufre una recalibración térmica que modifica la ganancia de la fuente controlada $E_1$ a $k = 1.5$ ($V_{\text{fuente}} = 1.5 V_x$).
    
    \begin{enumerate}
        \item \textbf{Cálculo Manuscrito:} Plantear la nueva matriz $\mathbf{A}$ y deducir los nuevos valores de $V_2$, $V_3$ y $A_v$.
        \item \textbf{Comprobación en Python:} Actualizar la matriz de coeficientes y resolver.
        \item \textbf{Validación en LTspice:} Correr \texttt{.tf} y registrar la nueva función de transferencia.
    \end{enumerate}
    
    \vspace{0.2cm}
    \begin{alertblock}{Resultados de Autocontrol}
        \begin{itemize}
            \item Ecuación de enlace: $V_2 - V_3 = 1.5 V_2 \implies \mathbf{-0.5 V_2 - V_3 = 0}$
            \item Voltajes resultantes: $\mathbf{V_2 = +7.273\text{ V}}, \quad \mathbf{V_3 = -3.636\text{ V}}$
            \item Ganancia de voltaje: $A_v = \frac{-3.636}{10} = \mathbf{-0.3636\text{ V/V}}$
            \item Error: $\varepsilon_r (\text{Teoría vs. LTspice}) = 0.00\%$
        \end{itemize}
    \end{alertblock}
\end{frame}

\end{document}